\documentclass[a4paper]{arrowhead}

\usepackage[yyyymmdd]{datetime}
\usepackage{etoolbox}
\usepackage[utf8]{inputenc}
\usepackage[T1]{fontenc}
\usepackage{multirow}
\usepackage{longtable}
\usepackage{tikz}
\usepackage{adjustbox}
\AtBeginEnvironment{longtable}{\small}
\BeforeBeginEnvironment{tikzpicture}{\begin{adjustbox}{max width=\textwidth}}
\AfterEndEnvironment{tikzpicture}{\end{adjustbox}}
\usetikzlibrary{positioning,shapes.multipart,arrows.meta}

\renewcommand{\dateseparator}{-}
\setlength{\emergencystretch}{3em}
\lstset{language=HTTP,breakatwhitespace=false,columns=fullflexible}

\newcommand{\fparam}[1]{\textit{\textcolor{ArrowheadBlue}{#1}}}

%% Special references
\newcommand{\mref}[2][]{{\textcolor{ArrowheadPurple}{\hyperref[sec:model:#2]{\ifstrempty{#1}{#2}{#1}}}}}
\newcommand{\oref}[1]{{\textcolor{ArrowheadBlue}{\hyperref[sec:operations:#1]{#1}}}}
\newcommand{\pdef}[1]{{\textcolor{ArrowheadGrey}{#1 \label{sec:model:primitives:#1} \label{sec:model:primitives:#1s}}}}
\newcommand{\pref}[1]{{\textcolor{ArrowheadGrey}{\hyperref[sec:model:primitives:#1]{#1}}}}

% as template_idd.tex; the input and output arguments are passed already formatted (\mref or \pref)
\newrobustcmd\osubsection[5]{
  \addtocounter{subsection}{1}
  \addcontentsline{toc}{subsection}{\protect\numberline{\thesubsection}operation \textcolor{ArrowheadBlue}{#1}}
  \subsection*{
    \thesubsection\quad
    #2 \textcolor{ArrowheadPurple}{#3} \\
    \small
    \hspace*{0.075\textwidth}\begin{minipage}{0.1\textwidth}
      \vspace*{1mm}
      Operation: \\
      \notblank{#4}{Input: \\}{}
      \notblank{#5}{Output: \\}{}
    \end{minipage}%
    \begin{minipage}{0.825\textwidth}
      \vspace*{1mm}
      \textcolor{ArrowheadBlue}{#1} \\
      \notblank{#4}{#4 \\}{}
      \notblank{#5}{#5 \\}{}
    \end{minipage}
  }
  \label{sec:operations:#1}
}
\newrobustcmd\msubsection[2]{
  \addtocounter{subsection}{1}
  \addcontentsline{toc}{subsection}{\protect\numberline{\thesubsection}#1 \textcolor{ArrowheadPurple}{#2}}
  \subsection*{\thesubsection\quad#1 \textcolor{ArrowheadPurple}{#2}}
  \label{sec:model:#2}
}
%%

\begin{document}
\sloppy

%% Arrowhead Document Properties
\ArrowheadTitle{serviceOrchestrationLockManagement MQTT/TLS/JSON}
\ArrowheadServiceID{serviceOrchestrationLockManagement}
\ArrowheadType{Interface Design Description}
\ArrowheadTypeShort{IDD}
\ArrowheadVersion{5.2.2}
\ArrowheadDate{2026-10-22}
\ArrowheadAuthor{Jerker Delsing}
\ArrowheadStatus{RELEASE CANDIDATE}
\ArrowheadContact{jerker.delsing@ltu.se}
\ArrowheadFooter{\href{www.arrowhead.eu}{www.arrowhead.eu}}
\ArrowheadSetup
%%

%% Front Page
\begin{center}
  \vspace*{1cm}
  \huge{\arrowtitle}

  \vspace*{0.2cm}
  \LARGE{\arrowtype}
  \vspace*{1cm}

  \Large{Service ID: \textit{"\arrowid"}}

  \vspace*{\fill}

  % Front Page Image

  \vspace*{1cm}
  \vspace*{\fill}

  % Front Page Abstract
  \begin{abstract}
    This document describes the MQTT protocol with TLS payload security and JSON payload encoding variant of the \texttt{service\allowbreak{}Orchestration\allowbreak{}Lock\allowbreak{}Management} service of the DynamicServiceOrchestration system.
  \end{abstract}

  \vspace*{1cm}
\end{center}

\newpage
%%

%% Table of Contents
\tableofcontents
\newpage
%%

\section{Overview}
\label{sec:overview}

This document describes the \texttt{service\allowbreak{}Orchestration\allowbreak{}Lock\allowbreak{}Management} service interface of the DynamicServiceOrchestration system over MQTT, which provides operator management of orchestration locks, the exclusive reservations of service instances.

It's implemented using protocol, encoding as stated in below table:
\begin{table}[ht!]
  \centering
  \begin{tabular}{|l|l|l|}
    \hline
    \rowcolor{gray!33} Profile type & Type & Version \\ \hline
    Transfer protocol & MQTT & 3.1.1 \cite{mqtt311} \\ \hline
    Data encryption & TLS & 1.3 \cite{rfc8446}, 1.2 \cite{rfc5246} \\ \hline
    Encoding & JSON & RFC 8259 \cite{rfc8259} \\ \hline
    Compression & N/A & - \\ \hline
    Semantics & N/A & - \\ \hline
    Ontology & N/A & - \\ \hline
  \end{tabular}
  \caption{Communication and semantics details used for the \texttt{service\allowbreak{}Orchestration\allowbreak{}Lock\allowbreak{}Management} service interface}
  \label{tab:comunication_semantics_profile}
\end{table}

This document provides the Interface Design Description IDD to the \textit{serviceOrchestrationLockManagement -- Service Description} document \cite{dso-serviceOrchestrationLockManagement-sd}.
The operations, their semantics and the data models are the same as in the HTTP/TLS/JSON IDD \cite{dso-serviceOrchestrationLockManagement-http-idd}; this document specifies the MQTT binding.

The MQTT interface is offered only when the system runs with \texttt{mqtt.api.enabled=true}; it is then registered with the interface template \texttt{generic\_\allowbreak{}mqtts} (\texttt{generic\_\allowbreak{}mqtt} without TLS), the broker address and port (\texttt{mqtt.\allowbreak{}broker.\allowbreak{}address}, \texttt{mqtt.\allowbreak{}broker.\allowbreak{}port}) and the base topic \texttt{arrowhead/\allowbreak{}serviceorchestration/\allowbreak{}orchestration/\allowbreak{}management/\allowbreak{}lock/\allowbreak{}}. In the Arrowhead-test-cloud-java deployment the MQTT interfaces of the core systems are disabled.

\newpage

\section{Service Operations}
\label{sec:operations}

The provider subscribes to one topic per operation, the base topic followed by the operation name:

\noindent\begin{longtable}{| p{5cm} | p{10.4cm} |} \hline
\rowcolor{gray!33} Operation & Request topic \\ \hline \endhead
\texttt{create} & \texttt{arrowhead/\allowbreak{}serviceorchestration/\allowbreak{}orchestration/\allowbreak{}management/\allowbreak{}lock/\allowbreak{}create} \\ \hline
\texttt{query} & \texttt{arrowhead/\allowbreak{}serviceorchestration/\allowbreak{}orchestration/\allowbreak{}management/\allowbreak{}lock/\allowbreak{}query} \\ \hline
\texttt{remove} & \texttt{arrowhead/\allowbreak{}serviceorchestration/\allowbreak{}orchestration/\allowbreak{}management/\allowbreak{}lock/\allowbreak{}remove} \\ \hline
\end{longtable}

A consumer publishes a request envelope (\mref{MqttRequestTemplate}) to the topic and receives the answer (\mref{MqttResponseTemplate}) on its own \texttt{responseTopic}. The \texttt{authentication} field carries the requester's identity according to the provider's authentication policy: \texttt{IDENTITY-TOKEN//<token>} (outsourced or internal policy), \texttt{SYSTEM//<name>} (declared policy) or the requester's X.509 certificate (certificate policy). The \texttt{status} of the response uses the HTTP status code numbers (200, 201, 204, 400, 401, 403, 404, 408, 423, 500, 503); on error the payload is an \mref{ErrorMessage}.

\begin{lstlisting}[caption={Request envelope.}]
{
  "traceId": "a1b2c3",
  "authentication": "IDENTITY-TOKEN//<token>",
  "responseTopic": "arrowhead/response/TemperatureDisplay/a1b2c3",
  "qosRequirement": 1,
  "params": {},
  "payload": {}
}
\end{lstlisting}

\begin{lstlisting}[caption={Response envelope.}]
{
  "status": 200,
  "traceId": "a1b2c3",
  "receiver": "TemperatureDisplay",
  "payload": {}
}
\end{lstlisting}

\osubsection{create}{PUBLISH}{\texttt{arrowhead/\allowbreak{}serviceorchestration/\allowbreak{}orchestration/\allowbreak{}management/\allowbreak{}lock/\allowbreak{}create}}{\mref[Orchestration\allowbreak{}Lock\allowbreak{}List\allowbreak{}Request]{OrchestrationLockListRequest}}{\mref[Orchestration\allowbreak{}Lock\allowbreak{}List\allowbreak{}Response]{OrchestrationLockListResponse}}

Creates locks on service instances for owners.

\begin{lstlisting}[caption={Payload of an \oref{create} request.}]
{
  "locks": [
    {
      "serviceInstanceId": "TemperatureSensor|temperatureReading|1.0.0",
      "owner": "TemperatureDisplay",
      "expiresAt": "2026-10-22T12:00:00Z"
    }
  ]
}
\end{lstlisting}

\osubsection{query}{PUBLISH}{\texttt{arrowhead/\allowbreak{}serviceorchestration/\allowbreak{}orchestration/\allowbreak{}management/\allowbreak{}lock/\allowbreak{}query}}{\mref[Orchestration\allowbreak{}Lock\allowbreak{}Query\allowbreak{}Request]{OrchestrationLockQueryRequest}}{\mref[Orchestration\allowbreak{}Lock\allowbreak{}List\allowbreak{}Response]{OrchestrationLockListResponse}}

Returns a page of locks matching the filter.

\begin{lstlisting}[caption={Payload of an \oref{query} request.}]
{
  "pagination": {
    "page": 0,
    "size": 10,
    "direction": "ASC",
    "sortField": "createdAt"
  },
  "ids": [
    1
  ],
  "orchestrationJobIds": [
    "3fa85f64-5717-4562-b3fc-2c963f66afa6"
  ],
  "serviceInstanceIds": [
    "TemperatureSensor|temperatureReading|1.0.0"
  ],
  "owners": [
    "TemperatureDisplay"
  ],
  "expiresBefore": "2026-10-22T12:00:00Z",
  "expiresAfter": "2026-10-22T12:00:00Z"
}
\end{lstlisting}

\osubsection{remove}{PUBLISH}{\texttt{arrowhead/\allowbreak{}serviceorchestration/\allowbreak{}orchestration/\allowbreak{}management/\allowbreak{}lock/\allowbreak{}remove}}{\pref{Array}$<$\allowbreak{}\pref{String}$>$}{}

Removes the locks of an owner on the given service instances.

\begin{lstlisting}[caption={Payload of an \oref{remove} request.}]
[
  "TemperatureDisplay"
]
\end{lstlisting}

\newpage

\section{Data Models}
\label{sec:model}

Here, all data objects that can be part of the service calls associated with this service are listed in alphabetic order.
Each \textit{struct} denotes a JSON \pref{Object} with the named fields; fields that are not needed may be omitted.

\msubsection{struct}{MqttRequestTemplate}

The request envelope.

\noindent\begin{tabularx}{\textwidth}{| p{4.25cm} | p{3.5cm} | X |} \hline
\rowcolor{gray!33} Object Field & Value Type & Description \\ \hline
"traceId" & \pref{String} & Identifier chosen by the requester, returned in the response. \\ \hline
"authentication" & \pref{String} & Identity of the requester (see Section \ref{sec:operations}). \\ \hline
"responseTopic" & \pref{String} & Topic on which the response is published. \\ \hline
"qosRequirement" & \pref{Number} & MQTT QoS of the response (0, 1 or 2). \\ \hline
"params" & \pref{Object}$<$\pref{String}$>$ & Optional parameters (the query parameters of the HTTP binding). \\ \hline
"payload" & \pref{Object} & The operation input. \\ \hline
\end{tabularx}

\msubsection{struct}{MqttResponseTemplate}

The response envelope.

\noindent\begin{tabularx}{\textwidth}{| p{4.25cm} | p{3.5cm} | X |} \hline
\rowcolor{gray!33} Object Field & Value Type & Description \\ \hline
"status" & \pref{Number} & Status code. \\ \hline
"traceId" & \pref{String} & The trace identifier of the request. \\ \hline
"receiver" & \pref{String} & Name of the requester. \\ \hline
"payload" & \pref{Object} & The operation output, or an \mref{ErrorMessage}. \\ \hline
\end{tabularx}

\msubsection{struct}{ErrorMessage}




\noindent\begin{tabularx}{\textwidth}{| >{\raggedright\arraybackslash}p{3.9cm} | >{\raggedright\arraybackslash\small}p{5.3cm} | X |} \hline
\rowcolor{gray!33} Object Field & Value Type & Description \\ \hline
"error\allowbreak{}Message" & \pref{String} &  \\ \hline
"error\allowbreak{}Code" & \pref{Number} &  \\ \hline
"exception\allowbreak{}Type" & \pref{ExceptionType} &  \\ \hline
"origin" & \pref{String} &  \\ \hline
\end{tabularx}



\msubsection{struct}{OrchestrationLockListRequest}

Locks to create.


\noindent\begin{tabularx}{\textwidth}{| >{\raggedright\arraybackslash}p{3.9cm} | >{\raggedright\arraybackslash\small}p{5.3cm} | X |} \hline
\rowcolor{gray!33} Object Field & Value Type & Description \\ \hline
"locks" & \pref{Array}$<$\allowbreak{}\mref{OrchestrationLockRequest}$>$ & Locks to create. \\ \hline
\end{tabularx}



\msubsection{struct}{OrchestrationLockListResponse}

List of orchestration locks.


\noindent\begin{tabularx}{\textwidth}{| >{\raggedright\arraybackslash}p{3.9cm} | >{\raggedright\arraybackslash\small}p{5.3cm} | X |} \hline
\rowcolor{gray!33} Object Field & Value Type & Description \\ \hline
"entries" & \pref{Array}$<$\allowbreak{}\mref[Orchestration\allowbreak{}Lock\allowbreak{}Response]{OrchestrationLockResponse}$>$ & Matching records of the current page. \\ \hline
"count" & \pref{Number} & Total number of matching records (all pages). \\ \hline
\end{tabularx}



\msubsection{struct}{OrchestrationLockQueryRequest}

Paged filter for orchestration locks.


\noindent\begin{tabularx}{\textwidth}{| >{\raggedright\arraybackslash}p{3.9cm} | >{\raggedright\arraybackslash\small}p{5.3cm} | X |} \hline
\rowcolor{gray!33} Object Field & Value Type & Description \\ \hline
"pagination" & \mref{Page} & Paging and sorting of the result. \\ \hline
"ids" & \pref{Array}$<$\allowbreak{}\pref{Number}$>$ & Lock identifiers to match. \\ \hline
"orchestration\allowbreak{}Job\allowbreak{}Ids" & \pref{Array}$<$\allowbreak{}\pref{String}$>$ & Orchestration job identifiers to match. \\ \hline
"service\allowbreak{}Instance\allowbreak{}Ids" & \pref{Array}$<$\allowbreak{}\pref{String}$>$ & Service instance identifiers to match. \\ \hline
"owners" & \pref{Array}$<$\allowbreak{}\pref{String}$>$ & Lock owner names to match. \\ \hline
"expires\allowbreak{}Before" & \pref{DateTime} & Upper bound of the expiry time. \\ \hline
"expires\allowbreak{}After" & \pref{DateTime} & Lower bound of the expiry time. \\ \hline
\end{tabularx}



\msubsection{struct}{OrchestrationLockRequest}

Exclusive reservation of a service instance.


\noindent\begin{tabularx}{\textwidth}{| >{\raggedright\arraybackslash}p{3.9cm} | >{\raggedright\arraybackslash\small}p{5.3cm} | X |} \hline
\rowcolor{gray!33} Object Field & Value Type & Description \\ \hline
"service\allowbreak{}Instance\allowbreak{}Id" & \pref{String} & Service instance identifier \texttt{\textless{}System\textgreater{}\textbar{}\allowbreak{}\textless{}service\allowbreak{}Definition\textgreater{}\textbar{}\allowbreak{}\textless{}version\textgreater{}}. \\ \hline
"owner" & \pref{String} & System for which the instance is locked. \\ \hline
"expires\allowbreak{}At" & \pref{DateTime} & UTC expiry of the lock. \\ \hline
\end{tabularx}



\msubsection{struct}{OrchestrationLockResponse}

A lock on a service instance.


\noindent\begin{tabularx}{\textwidth}{| >{\raggedright\arraybackslash}p{3.9cm} | >{\raggedright\arraybackslash\small}p{5.3cm} | X |} \hline
\rowcolor{gray!33} Object Field & Value Type & Description \\ \hline
"id" & \pref{Number} & Identifier of the lock. \\ \hline
"orchestration\allowbreak{}Job\allowbreak{}Id" & \pref{String} & Orchestration job that created the lock. \\ \hline
"service\allowbreak{}Instance\allowbreak{}Id" & \pref{String} & Service instance identifier \texttt{\textless{}System\textgreater{}\textbar{}\allowbreak{}\textless{}service\allowbreak{}Definition\textgreater{}\textbar{}\allowbreak{}\textless{}version\textgreater{}}. \\ \hline
"owner" & \pref{String} & Name of the system owning the lock. \\ \hline
"expires\allowbreak{}At" & \pref{DateTime} & UTC expiry of the lock. \\ \hline
"temporary" & \pref{Boolean} & True for a lock held only during an orchestration. \\ \hline
\end{tabularx}



\msubsection{struct}{Page}

Paging and sorting; all fields optional.


\noindent\begin{tabularx}{\textwidth}{| >{\raggedright\arraybackslash}p{3.9cm} | >{\raggedright\arraybackslash\small}p{5.3cm} | X |} \hline
\rowcolor{gray!33} Object Field & Value Type & Description \\ \hline
"page" & \pref{Number} & Zero-based page number. \\ \hline
"size" & \pref{Number} & Number of records per page (at most \texttt{max.\allowbreak{}page.\allowbreak{}size}, default 1000). \\ \hline
"direction" & \pref{String} & Sort direction, \texttt{ASC} or \texttt{DESC}. \\ \hline
"sort\allowbreak{}Field" & \pref{String} & Field to sort by. \\ \hline
\end{tabularx}



\subsection{Primitives}
\label{sec:model:primitives}

As all messages are encoded using the JSON format \cite{rfc8259}, the following primitive constructs, part of that standard, become available.

\noindent\begin{tabularx}{\textwidth}{| p{4.6cm} | X |} \hline
\rowcolor{gray!33} Type & Description \\ \hline
\pdef{Array} & A JSON Array \cite{rfc8259}; \pref{Array}$<$A$>$ has elements of type A. \\ \hline
\pdef{Boolean} & JSON Boolean, \texttt{true} or \texttt{false}. \\ \hline
\pdef{DateTime} & JSON String with an ISO 8601 UTC timestamp, e.g. \texttt{2026-10-22T12:00:00Z}. \\ \hline
\pdef{ExceptionType} & One of \texttt{ARROWHEAD}, \texttt{INVALID\_\allowbreak{}PARAMETER}, \texttt{AUTH}, \texttt{FORBIDDEN}, \texttt{DATA\_\allowbreak{}NOT\_\allowbreak{}FOUND}, \texttt{TIMEOUT}, \texttt{LOCKED}, \texttt{INTERNAL\_\allowbreak{}SERVER\_\allowbreak{}ERROR}, \texttt{EXTERNAL\_\allowbreak{}SERVER\_\allowbreak{}ERROR} (JSON String). \\ \hline
\pdef{Number} & JSON Number, an integer. \\ \hline
\pdef{Object} & JSON Object; \pref{Object}$<$A$>$ has string keys and values of type A. \\ \hline
\pdef{String} & JSON String. \\ \hline
\end{tabularx}

\newpage

\bibliographystyle{IEEEtran}
\bibliography{bibliography}

\newpage

\section{Revision History}
\subsection{Amendments}

\noindent\begin{tabularx}{\textwidth}{| p{1cm} | p{3cm} | p{2cm} | X | p{4cm} |} \hline
\rowcolor{gray!33} No. & Date & Version & Subject of Amendments & Author \\ \hline

1 & 2026-10-22 & \arrowversion & Document created from the Arrowhead LaTeX templates for the 5.2.2 Java Spring implementation, including the local extensions of this workspace. & Jerker Delsing \\ \hline
\end{tabularx}

\subsection{Quality Assurance}

\noindent\begin{tabularx}{\textwidth}{| p{1cm} | p{3cm} | p{2cm} | X |} \hline
\rowcolor{gray!33} No. & Date & Version & Approved by \\ \hline

1 & & \arrowversion  &  \\ \hline

\end{tabularx}

\end{document}
